\begin{figure}[htbp]
\centering
\begin{tikzpicture}[fd,x=1cm,y=1cm,
  fd2 dot/.style={circle,fill=#1,draw=#1!70!black,line width=0.4pt,inner sep=0pt,minimum size=0.176cm},
  fd2 hol/.style={circle,fill=white,draw=#1,line width=0.85pt,inner sep=0pt,minimum size=0.166cm},
  fd2 row/.style={font=\footnotesize,anchor=west,inner sep=0pt},
  fd2 ids/.style={font=\tiny,anchor=west,inner sep=0pt,text=gray!45!black},
  fd2 fn/.style={font=\tiny,anchor=west,inner sep=0.5pt}]
  \fill[gray!7] (0,-1.440) rectangle (15.950,-0.720);
  \fill[gray!7] (0,-2.880) rectangle (15.950,-2.160);
  \fill[gray!7] (0,-4.320) rectangle (15.950,-3.600);
  \fill[gray!7] (0,-5.760) rectangle (15.950,-5.040);
  \fill[gray!7] (0,-7.200) rectangle (15.950,-6.480);
  \fill[gray!7] (0,-8.640) rectangle (15.950,-7.920);
  \fill[gray!7] (0,-10.080) rectangle (15.950,-9.360);
  \fill[gray!7] (0,-11.520) rectangle (15.950,-10.800);
  \fill[gray!7] (0,-12.960) rectangle (15.950,-12.240);
  \fill[gray!7] (0,-14.400) rectangle (15.950,-13.680);
  \node[inner sep=0pt] at (5.600,0.420) {\tagR};
  \node[inner sep=0pt] at (7.900,0.420) {\tagA};
  \node[inner sep=0pt] at (10.200,0.420) {\tagP};
  \node[inner sep=0pt] at (12.500,0.420) {\tagC};
  \node[inner sep=0pt] at (14.800,0.420) {\tagV};
  \node[font=\footnotesize\bfseries,anchor=west,inner sep=0pt] at (0,0.420) {过程类};
  \draw[gray!60,line width=0.6pt] (0,0.000) -- (15.950,0.000);
  \draw[gray!55,line width=0.6pt] (0,-5.040) -- (15.950,-5.040);
  \draw[gray!55,line width=0.6pt] (0,-9.360) -- (15.950,-9.360);
  \draw[gray!55,line width=0.6pt] (0,-12.240) -- (15.950,-12.240);
  \draw[gray!60,line width=0.6pt] (0,-14.400) -- (15.950,-14.400);
  \draw[gray!35,line width=0.4pt] (4.450,0.000) -- (4.450,-14.400);
  \draw[gray!35,line width=0.4pt] (6.750,0.000) -- (6.750,-14.400);
  \draw[gray!35,line width=0.4pt] (9.050,0.000) -- (9.050,-14.400);
  \draw[gray!35,line width=0.4pt] (11.350,0.000) -- (11.350,-14.400);
  \draw[gray!35,line width=0.4pt] (13.650,0.000) -- (13.650,-14.400);
  \draw[gray!35,line width=0.4pt] (15.950,0.000) -- (15.950,-14.400);
  \node[fd2 row] at (0.06,-0.250) {$e^+e^-$ 总截面};
  \node[fd2 ids] at (0.06,-0.560) {II.1--7};
  \node[fd2 dot=OliveGreen] at (4.660,-0.360) {};  % II.1
  \node[fd2 dot=OliveGreen] at (4.888,-0.360) {};  % II.2
  \node[fd2 dot=OliveGreen] at (5.116,-0.360) {};  % II.4
  \node[fd2 dot=RoyalBlue] at (6.960,-0.360) {};  % II.3
  \node[fd2 dot=RoyalBlue] at (7.188,-0.360) {};  % II.5
  \node[fd2 dot=Plum] at (11.560,-0.360) {};  % II.6
  \node[fd2 dot=Plum] at (11.788,-0.360) {};  % II.7
  \node[fd2 row] at (0.06,-0.970) {DIS 矩、OPE 与极化求和规则};
  \node[fd2 ids] at (0.06,-1.280) {II.8--15,\,17--19,\,28,\,29；III.21--24};
  \node[fd2 dot=OliveGreen] at (4.660,-1.080) {};  % II.17
  \node[fd2 dot=OliveGreen] at (4.888,-1.080) {};  % II.18
  \node[fd2 dot=OliveGreen] at (5.116,-1.080) {};  % II.28
  \node[fd2 dot=OliveGreen] at (5.344,-1.080) {};  % III.21
  \node[fd2 hol=OliveGreen] at (5.632,-1.080) {};  % II.12
  \node[fd2 hol=OliveGreen] at (5.860,-1.080) {};  % II.13
  \node[fd2 hol=OliveGreen] at (6.088,-1.080) {};  % II.14
  \node[fd2 hol=OliveGreen] at (6.316,-1.080) {};  % II.15
  \node[fd2 fn] at (6.434,-0.990) {(a)};
  \node[fd2 dot=RoyalBlue] at (6.960,-1.080) {};  % II.8
  \node[fd2 dot=RoyalBlue] at (7.188,-1.080) {};  % II.9
  \node[fd2 dot=RoyalBlue] at (7.416,-1.080) {};  % II.10
  \node[fd2 dot=RoyalBlue] at (7.644,-1.080) {};  % II.11
  \node[fd2 dot=RoyalBlue] at (7.872,-1.080) {};  % II.19
  \node[fd2 hol=RoyalBlue] at (8.160,-1.080) {};  % II.29
  \node[fd2 dot=Plum] at (11.560,-1.080) {};  % III.23
  \node[fd2 dot=BrickRed] at (13.860,-1.080) {};  % III.22
  \node[fd2 dot=BrickRed] at (14.088,-1.080) {};  % III.24
  \node[fd2 row] at (0.06,-1.690) {DIS $x$ 空间};
  \node[fd2 ids] at (0.06,-2.000) {II.16,\,25；III.10--15；VII.30};
  \node[fd2 dot=OliveGreen] at (4.660,-1.800) {};  % II.25
  \node[fd2 dot=OliveGreen] at (4.888,-1.800) {};  % III.14
  \node[fd2 hol=OliveGreen] at (5.176,-1.800) {};  % II.16
  \node[fd2 fn] at (5.294,-1.710) {(a)};
  \node[fd2 dot=RoyalBlue] at (6.960,-1.800) {};  % III.13
  \node[fd2 dot=RoyalBlue] at (7.188,-1.800) {};  % III.15
  \node[fd2 dot=RoyalBlue] at (7.416,-1.800) {};  % VII.30
  \node[fd2 hol=RoyalBlue] at (7.704,-1.800) {};  % III.10
  \node[fd2 hol=RoyalBlue] at (7.932,-1.800) {};  % III.11
  \node[fd2 hol=Plum] at (11.560,-1.800) {};  % III.12
  \node[fd2 fn] at (11.678,-1.710) {(c)};
  \node[fd2 row] at (0.06,-2.410) {SIA、SIDIS 与光生};
  \node[fd2 ids] at (0.06,-2.720) {III.16--19};
  \node[fd2 dot=RoyalBlue] at (6.960,-2.520) {};  % III.16
  \node[fd2 dot=RoyalBlue] at (7.188,-2.520) {};  % III.17
  \node[fd2 dot=RoyalBlue] at (7.416,-2.520) {};  % III.18
  \node[fd2 dot=Plum] at (11.560,-2.520) {};  % III.19
  \node[fd2 row] at (0.06,-3.130) {DY 与色单态末态};
  \node[fd2 ids] at (0.06,-3.440) {III.25--30,\,33,\,34};
  \node[fd2 dot=OliveGreen] at (4.660,-3.240) {};  % III.27
  \node[fd2 dot=RoyalBlue] at (6.960,-3.240) {};  % III.26
  \node[fd2 dot=RoyalBlue] at (7.188,-3.240) {};  % III.29
  \node[fd2 dot=RoyalBlue] at (7.416,-3.240) {};  % III.30
  \node[fd2 dot=RoyalBlue] at (7.644,-3.240) {};  % III.33
  \node[fd2 hol=RoyalBlue] at (7.932,-3.240) {};  % III.25
  \node[fd2 dot=BurntOrange] at (9.260,-3.240) {};  % III.28
  \node[fd2 dot=BrickRed] at (13.860,-3.240) {};  % III.34
  \node[fd2 row] at (0.06,-3.850) {$pp\to$ 喷注/强子（有色末态）};
  \node[fd2 ids] at (0.06,-4.160) {III.31,\,32；IV.24--26};
  \node[fd2 dot=Plum] at (11.560,-3.960) {};  % III.31
  \node[fd2 dot=Plum] at (11.788,-3.960) {};  % III.32
  \node[fd2 dot=Plum] at (12.016,-3.960) {};  % IV.25
  \node[fd2 dot=BrickRed] at (13.860,-3.960) {};  % IV.24
  \node[fd2 dot=BrickRed] at (14.088,-3.960) {};  % IV.26
  \node[fd2 row] at (0.06,-4.570) {阈值重新因子化};
  \node[fd2 ids] at (0.06,-4.880) {III.35--37};
  \node[fd2 dot=RoyalBlue] at (6.960,-4.680) {};  % III.35
  \node[fd2 dot=RoyalBlue] at (7.188,-4.680) {};  % III.36
  \node[fd2 dot=BurntOrange] at (9.260,-4.680) {};  % III.37
  \node[fd2 row] at (0.06,-5.290) {TMD：DY、SIDIS、$e^+e^-$};
  \node[fd2 ids] at (0.06,-5.600) {IV.1--9；VII.48};
  \node[fd2 dot=OliveGreen] at (4.660,-5.400) {};  % IV.2
  \node[fd2 dot=OliveGreen] at (4.888,-5.400) {};  % IV.6
  \node[fd2 dot=RoyalBlue] at (6.960,-5.400) {};  % IV.3
  \node[fd2 dot=RoyalBlue] at (7.188,-5.400) {};  % IV.4
  \node[fd2 dot=RoyalBlue] at (7.416,-5.400) {};  % IV.5
  \node[fd2 dot=RoyalBlue] at (7.644,-5.400) {};  % IV.7
  \node[fd2 dot=RoyalBlue] at (7.872,-5.400) {};  % IV.8
  \node[fd2 hol=RoyalBlue] at (8.160,-5.400) {};  % IV.1
  \node[fd2 dot=BurntOrange] at (9.260,-5.400) {};  % IV.9
  \node[fd2 dot=Plum] at (11.560,-5.400) {};  % VII.48
  \node[fd2 row] at (0.06,-6.010) {TMD：$h_1h_2\to H_1H_2X$};
  \node[fd2 ids] at (0.06,-6.320) {IV.22,\,23};
  \node[fd2 hol=RoyalBlue] at (6.960,-6.120) {};  % IV.22
  \node[fd2 dot=BrickRed] at (13.860,-6.120) {};  % IV.23
  \node[fd2 row] at (0.06,-6.730) {SCET、事件形状与喷注};
  \node[fd2 ids] at (0.06,-7.040) {IV.10--17；VII.23--25};
  \node[fd2 dot=OliveGreen] at (4.660,-6.840) {};  % IV.10
  \node[fd2 dot=RoyalBlue] at (6.960,-6.840) {};  % IV.11
  \node[fd2 dot=RoyalBlue] at (7.188,-6.840) {};  % IV.12
  \node[fd2 dot=RoyalBlue] at (7.416,-6.840) {};  % IV.13
  \node[fd2 dot=RoyalBlue] at (7.644,-6.840) {};  % IV.15
  \node[fd2 dot=RoyalBlue] at (7.872,-6.840) {};  % VII.23
  \node[fd2 dot=BurntOrange] at (9.260,-6.840) {};  % IV.14
  \node[fd2 dot=BurntOrange] at (9.488,-6.840) {};  % IV.16
  \node[fd2 dot=BurntOrange] at (9.716,-6.840) {};  % VII.24
  \node[fd2 dot=Plum] at (11.560,-6.840) {};  % IV.17
  \node[fd2 dot=BrickRed] at (13.860,-6.840) {};  % VII.25
  \node[fd2 row] at (0.06,-7.450) {小 $x$};
  \node[fd2 ids] at (0.06,-7.760) {IV.18--21；VII.31};
  \node[fd2 dot=RoyalBlue] at (6.960,-7.560) {};  % IV.18
  \node[fd2 dot=RoyalBlue] at (7.188,-7.560) {};  % IV.20
  \node[fd2 dot=BurntOrange] at (9.260,-7.560) {};  % IV.19
  \node[fd2 dot=Plum] at (11.560,-7.560) {};  % IV.21
  \node[fd2 dot=BrickRed] at (13.860,-7.560) {};  % VII.31
  \node[fd2 row] at (0.06,-8.170) {衍射};
  \node[fd2 ids] at (0.06,-8.480) {III.20；IV.28,\,29};
  \node[fd2 dot=RoyalBlue] at (6.960,-8.280) {};  % IV.28
  \node[fd2 dot=BrickRed] at (13.860,-8.280) {};  % III.20
  \node[fd2 dot=BrickRed] at (14.088,-8.280) {};  % IV.29
  \node[fd2 row] at (0.06,-8.890) {双部分子散射与多重相互作用};
  \node[fd2 ids] at (0.06,-9.200) {IV.27,\,30--32；VII.22,\,46};
  \node[fd2 dot=RoyalBlue] at (6.960,-9.000) {};  % IV.30
  \node[fd2 dot=RoyalBlue] at (7.188,-9.000) {};  % VII.22
  \node[fd2 dot=Plum] at (11.560,-9.000) {};  % IV.31
  \node[fd2 dot=Plum] at (11.788,-9.000) {};  % IV.32
  \node[fd2 dot=Plum] at (12.016,-9.000) {};  % VII.46
  \node[fd2 dot=BrickRed] at (13.860,-9.000) {};  % IV.27
  \node[fd2 row] at (0.06,-9.610) {硬遍举过程};
  \node[fd2 ids] at (0.06,-9.920) {V.1--8};
  \node[fd2 dot=RoyalBlue] at (6.960,-9.720) {};  % V.2
  \node[fd2 dot=RoyalBlue] at (7.188,-9.720) {};  % V.3
  \node[fd2 dot=RoyalBlue] at (7.416,-9.720) {};  % V.5
  \node[fd2 dot=RoyalBlue] at (7.644,-9.720) {};  % V.6
  \node[fd2 hol=RoyalBlue] at (7.932,-9.720) {};  % V.1
  \node[fd2 dot=BurntOrange] at (9.260,-9.720) {};  % V.4
  \node[fd2 dot=BurntOrange] at (9.488,-9.720) {};  % V.8
  \node[fd2 dot=BrickRed] at (13.860,-9.720) {};  % V.7
  \node[fd2 row] at (0.06,-10.330) {$B$ 介子衰变};
  \node[fd2 ids] at (0.06,-10.640) {V.9--13};
  \node[fd2 dot=RoyalBlue] at (6.960,-10.440) {};  % V.9
  \node[fd2 dot=RoyalBlue] at (7.188,-10.440) {};  % V.10
  \node[fd2 dot=RoyalBlue] at (7.416,-10.440) {};  % V.13
  \node[fd2 dot=BurntOrange] at (9.260,-10.440) {};  % V.11
  \node[fd2 dot=BrickRed] at (13.860,-10.440) {};  % V.12
  \node[fd2 row] at (0.06,-11.050) {重夸克偶素（NRQCD）};
  \node[fd2 ids] at (0.06,-11.360) {V.14--16};
  \node[fd2 dot=RoyalBlue] at (6.960,-11.160) {};  % V.14
  \node[fd2 dot=RoyalBlue] at (7.188,-11.160) {};  % V.16
  \node[fd2 dot=Plum] at (11.560,-11.160) {};  % V.15
  \node[fd2 row] at (0.06,-11.770) {幂次修正、NLP 与重整化子};
  \node[fd2 ids] at (0.06,-12.080) {V.17--24,\,26；VII.26,\,47};
  \node[fd2 dot=OliveGreen] at (4.660,-11.880) {};  % V.21
  \node[fd2 dot=RoyalBlue] at (6.960,-11.880) {};  % V.17
  \node[fd2 dot=RoyalBlue] at (7.188,-11.880) {};  % VII.26
  \node[fd2 dot=BurntOrange] at (9.260,-11.880) {};  % V.18
  \node[fd2 dot=BurntOrange] at (9.488,-11.880) {};  % V.20
  \node[fd2 dot=BurntOrange] at (9.716,-11.880) {};  % V.22
  \node[fd2 dot=BurntOrange] at (9.944,-11.880) {};  % V.23
  \node[fd2 dot=BurntOrange] at (10.172,-11.880) {};  % V.26
  \node[fd2 dot=Plum] at (11.560,-11.880) {};  % V.19
  \node[fd2 dot=Plum] at (11.788,-11.880) {};  % V.24
  \node[fd2 dot=Plum] at (12.016,-11.880) {};  % VII.47
  \node[fd2 row] at (0.06,-12.490) {格点与 LaMET（共线）};
  \node[fd2 ids] at (0.06,-12.800) {VII.1--9,\,16,\,17,\,33--38,\,40--42};
  \node[fd2 dot=OliveGreen] at (4.660,-12.600) {};  % VII.1
  \node[fd2 dot=OliveGreen] at (4.888,-12.600) {};  % VII.6
  \node[fd2 dot=OliveGreen] at (5.116,-12.600) {};  % VII.7
  \node[fd2 dot=RoyalBlue] at (6.960,-12.600) {};  % VII.2
  \node[fd2 dot=RoyalBlue] at (7.188,-12.600) {};  % VII.3
  \node[fd2 dot=RoyalBlue] at (7.416,-12.600) {};  % VII.16
  \node[fd2 dot=RoyalBlue] at (7.644,-12.600) {};  % VII.34
  \node[fd2 dot=RoyalBlue] at (7.872,-12.600) {};  % VII.35
  \node[fd2 dot=RoyalBlue] at (8.100,-12.600) {};  % VII.38
  \node[fd2 hol=RoyalBlue] at (8.388,-12.600) {};  % VII.5
  \node[fd2 hol=RoyalBlue] at (8.616,-12.600) {};  % VII.33
  \node[fd2 dot=BurntOrange] at (9.260,-12.600) {};  % VII.17
  \node[fd2 dot=Plum] at (11.560,-12.600) {};  % VII.4
  \node[fd2 dot=Plum] at (11.788,-12.600) {};  % VII.36
  \node[fd2 dot=Plum] at (12.016,-12.600) {};  % VII.40
  \node[fd2 dot=Plum] at (12.244,-12.600) {};  % VII.41
  \node[fd2 dot=Plum] at (12.472,-12.600) {};  % VII.42
  \node[fd2 hol=Plum] at (12.760,-12.600) {};  % VII.8
  \node[fd2 hol=Plum] at (12.988,-12.600) {};  % VII.9
  \node[fd2 fn] at (13.106,-12.510) {(b)};
  \node[fd2 dot=BrickRed] at (13.860,-12.600) {};  % VII.37
  \node[fd2 row] at (0.06,-13.210) {准 TMD 与软函数};
  \node[fd2 ids] at (0.06,-13.520) {VII.10--15,\,39};
  \node[fd2 dot=RoyalBlue] at (6.960,-13.320) {};  % VII.10
  \node[fd2 dot=RoyalBlue] at (7.188,-13.320) {};  % VII.11
  \node[fd2 dot=RoyalBlue] at (7.416,-13.320) {};  % VII.12
  \node[fd2 dot=RoyalBlue] at (7.644,-13.320) {};  % VII.13
  \node[fd2 dot=RoyalBlue] at (7.872,-13.320) {};  % VII.14
  \node[fd2 dot=RoyalBlue] at (8.100,-13.320) {};  % VII.15
  \node[fd2 dot=Plum] at (11.560,-13.320) {};  % VII.39
  \node[fd2 row] at (0.06,-13.930) {Landau 理论、区域与 OPE};
  \node[fd2 ids] at (0.06,-14.240) {I.10--12,\,22,\,23；V.25；VII.18--21,\,43--45};
  \node[fd2 dot=OliveGreen] at (4.660,-14.040) {};  % I.10
  \node[fd2 dot=OliveGreen] at (4.888,-14.040) {};  % I.11
  \node[fd2 dot=OliveGreen] at (5.116,-14.040) {};  % I.12
  \node[fd2 dot=OliveGreen] at (5.344,-14.040) {};  % I.22
  \node[fd2 dot=OliveGreen] at (5.572,-14.040) {};  % V.25
  \node[fd2 dot=RoyalBlue] at (6.960,-14.040) {};  % VII.18
  \node[fd2 hol=RoyalBlue] at (7.248,-14.040) {};  % VII.19
  \node[fd2 dot=BurntOrange] at (9.260,-14.040) {};  % I.23
  \node[fd2 dot=Plum] at (11.560,-14.040) {};  % VII.20
  \node[fd2 dot=Plum] at (11.788,-14.040) {};  % VII.45
  \node[fd2 dot=BrickRed] at (13.860,-14.040) {};  % VII.21
  \node[fd2 dot=BrickRed] at (14.088,-14.040) {};  % VII.43
  \node[fd2 dot=BrickRed] at (14.316,-14.040) {};  % VII.44
  \node[fd2 dot=gray,anchor=center] at (0.15,-14.820) {};
  \node[font=\footnotesize,anchor=west,inner sep=0pt] at (0.32,-14.820) {表~\ref{tab:summary-master} 中的一条陈述，画在其规范标签所在列（$R_{\mathrm{PT}}$、$R_{\mathrm{def}}$、$R_{\mathrm{OPE}}$ 均在 [R] 列）；颜色同标签};
  \node[fd2 hol=darkgray,anchor=center] at (0.15,-15.240) {};
  \node[font=\footnotesize,anchor=west,inner sep=0pt] at (0.32,-15.240) {条件定理 $[\mathrm R,R_{\mathrm{cond}}]$，画在其内容等级所在列（全局假设 \texttt{asm:intro-*} 不降低内容等级）};
  \node[font=\footnotesize,anchor=north west,inner sep=0pt,align=left,text width=15.85cm] at (0.06,-15.590) {[A] 列中的空心圆：部分子层次的条件定理，前提为 [A] 级技术假设（\texttt{asm:appM-*}、\texttt{asm:appN-*}、\texttt{asm:appO-*}、\cref{asm:lamet-ope}）或 [A] 级定理。(a)~前提只含全局非微扰假设，II.14、II.15 的瓶颈为 \cref{asm:intro-np-ope,asm:intro-growth}；\ (b)~强子层次，瓶颈为 [C] 级假设：VII.8 为 \cref{asm:lamet-uniform,asm:lamet-twist4,asm:lamet-large-z}，VII.9 为 \cref{asm:lamet-gap}；\ (c)~部分子层次的正规子路线，内容等级经由 \cref{asm:appN-ir-regulator} 为 [C]。};
\end{tikzpicture}
\caption{各类过程中因子化陈述的地位一览（与表~\ref{tab:summary-master} 一致）。每个圆点代表该表的一行：实心圆画在其规范标签所在的列，空心圆为条件定理 $[\mathrm R,R_{\mathrm{cond}}]$，画在其内容等级所在的列，[C] 列的空心圆是内容等级为 [C] 的条件定理（VII.8、VII.9 为强子层次的陈述），图下注明其瓶颈假设；左栏小字为归入该类的表行编号。基础、方法与部分子分布性质的陈述（区块 I 的大部分、II.20--II.24、II.26、II.27、III.1--III.9、III.38）以及 VII.27--VII.29、VII.32 不属于某一过程类，未画出。同一过程类往往同时含有不同等级的陈述：标签属于陈述，而不属于过程。}
\label{fig:summary-status-map}
\end{figure}
